\documentclass[12pt]{article}
\usepackage[utf8]{inputenc}
\usepackage[T1]{fontenc}
\usepackage{hyperref}
\title{Northwestern Tamang Language Resources for K-12 Curriculum Development}
\author{Compiled by Claude Sonnet 5 (claude-sonnet-5) for Andr\'as Kornai}
\date{}
\begin{document}
\maketitle

\section{Introduction}
Northwestern Tamang (ISO 639-3 \texttt{tmk}) is a Tamang language of Nepal. The Asia Harvest profile gives 102,300 speakers (2020), concentrated in Nuwakot District. Most Tamang materials found are for Western Tamang, a different variety, so they are not direct resources for \texttt{tmk}.

\section{Primary Educational Resources}
No \texttt{tmk}-specific textbook was found. Omniglot says Western and Northwestern Tamang are taught in primary schools: \url{https://omniglot.com/writing/tamang.htm}. A study on Tamang mother-tongue-based multilingual education in Nepal reports a shortage of teaching and learning materials: \url{https://www.nepjol.info/index.php/mef/article/view/97701}

\section{Linguistic and Reference Materials}
\begin{itemize}
\item Western Tamang--Nepali dictionary on Webonary (a different variety): \url{https://www.webonary.org/wtamang/overview/alphabet/}
\item Wikipedia, Western Tamang: \url{https://en.wikipedia.org/wiki/Western_Tamang}
\end{itemize}

\section{Religious and Community Resources}
None found for this variety. The Asia Harvest people-group profile is at \url{https://people-groups.asiaharvest.org/Nepal/TamangNorthwestern.pdf}.

\section{Digital and Audio Resources}
Webonary hosts the Western Tamang dictionary (see above). No audio resources were found.

\section{Curriculum Development Context}
Tamang mother-tongue education exists in Nepal, but the materials found are for Western Tamang and the sources report a shortage of materials. Checking whether the national Tamang curriculum materials from the Curriculum Development Centre suit the Northwestern variety is recommended.
\end{document}